% Generated draft: requires mathematical and source review.
\documentclass[11pt]{article}
\usepackage[margin=25mm]{geometry}
\usepackage{fontspec}
\setmainfont{Latin Modern Roman}
\usepackage{amsmath,amssymb,mathtools,unicode-math}
\setmathfont{Latin Modern Math}
\usepackage{enumitem,needspace,xurl,hyperref}
\setlength{\parindent}{0pt}
\setlength{\parskip}{5pt}
\setlength{\emergencystretch}{3em}
\newcommand{\N}{\mathbb{N}}
\newcommand{\Z}{\mathbb{Z}}
\newcommand{\Q}{\mathbb{Q}}
\newcommand{\R}{\mathbb{R}}
\newcommand{\C}{\mathbb{C}}
\newcommand{\F}{\mathbb{F}}
\newcommand{\A}{\mathbb{A}}
\newcommand{\PP}{\mathbb{P}}
\newcommand{\E}{\mathbb{E}}
\newcommand{\eps}{\varepsilon}
\newcommand{\dd}{\,\mathrm{d}}
\begin{document}
\section*{UnsolvedMath ID: 700350}
\textbf{Source problem code:} catalogue entry
\subsection{Definitions and mathematical statement}
Let \(G\) be a group, let \(\Z G\) denote its integral group ring, and regard \(\Z\) as the trivial left \(\Z G\)-module, so that every element of \(G\) acts as the identity. The group \(G\) is of type \((FP)_{\infty}\) if there is an exact sequence of \(\Z G\)-modules
\[
\cdots \longrightarrow P_2 \longrightarrow P_1 \longrightarrow P_0 \longrightarrow \Z \longrightarrow 0
\]
in which \(P_n\) is finitely generated and projective for every integer \(n\geq 0\). The resolution need not have finite length. A group is periodic if each of its elements has finite order; no uniform bound on these orders is assumed.

\textbf{Problem (R. Bieri, 1978).} Is every periodic group of type \((FP)_{\infty}\) finite? Equivalently, is the following implication valid for every group \(G\)?
\[
\left[\left(\forall g\in G\;\exists m\in\Z_{\geq 1}:g^m=1_G\right)\ \text{and}\ G\text{ is of type }(FP)_{\infty}\right]
\quad\Longrightarrow\quad |G|<\infty.
\]

The supplied source describes closure properties of groups of type \((FP)_{\infty}\) involving extensions and amalgamated products and cites Bieri's 1976 monograph, but does not state the hypotheses of those closure results. They are background rather than additional assumptions of this problem. The source also relates the problem to the question whether an infinite periodic group can have a finite presentation; it does not assert that the two questions are equivalent.
\subsection{Short English statement}
Must a group in which every element has finite order be finite if its trivial integral module admits a projective resolution with finitely generated modules in every degree? The supplied record labels this question open; that status has not been independently verified.
\subsection{Sources}
\begin{itemize}
\item Local UnsolvedMath registry, permanent ID \(700350\), exact source problem code \texttt{catalogue entry}, title \emph{Kourovka Notebook Problem 6.3}. The supplied record attributes the problem to R. Bieri, with proposed year 1978, and labels it open.
\item \emph{The Kourovka Notebook (Unsolved Problems in Group Theory)}, original problem 6.3, 6th issue (1978); supplied edition reference: 21st issue edition (2026), \texttt{arXiv:1401.0300v48}. The supplied collection description names E. I. Khukhro and V. D. Mazurov as editors. These bibliographic details are supplied attribution and have not been independently checked.
\item R. Bieri, \emph{Homological dimension of discrete groups}, Queen Mary College Math. Notes, London, 1976. Supplied reference for the background discussion of homological finiteness and closure properties; no page or theorem number is supplied.
\end{itemize}
\end{document}
